\documentclass[11pt]{beamer}

\usepackage[T1]{fontenc}
\usepackage{tikz}
\usepackage[quiet]{mathspec}
\usepackage{ifthen}

\usetikzlibrary{matrix,arrows.meta}

% The base16 colour scheme?
\definecolor{sbase03}{HTML}{002B36}
\definecolor{sbase02}{HTML}{073642}
\definecolor{sbase01}{HTML}{586E75}
\definecolor{sbase00}{HTML}{657B83}
\definecolor{sbase0}{HTML}{839496}
\definecolor{sbase1}{HTML}{93A1A1}
\definecolor{sbase2}{HTML}{EEE8D5}
\definecolor{sbase3}{HTML}{FDF6E3}
\definecolor{syellow}{HTML}{B58900}
\definecolor{sorange}{HTML}{CB4B16}
\definecolor{sred}{HTML}{DC322F}
\definecolor{smagenta}{HTML}{D33682}
\definecolor{sviolet}{HTML}{6C71C4}
\definecolor{sblue}{HTML}{268BD2}
\definecolor{scyan}{HTML}{2AA198}
\definecolor{sgreen}{HTML}{859900}

\definecolor{Tropiteal}{RGB}{0,168,198}
\definecolor{TealDrop}{RGB}{64,192,203}
\definecolor{WhiteTrash}{RGB}{249,242,231}
\definecolor{AtomicBikini}{RGB}{174,226,57}
\definecolor{FeebleWeek}{RGB}{143,190,0}
\definecolor{ICantExpress}{RGB}{28,20,13}
\definecolor{Marty}{RGB}{250,42,0}

\colorlet{ColourBase}{Tropiteal}
\colorlet{ColourHl1}{Marty}
\colorlet{ColourHl2}{FeebleWeek}
\colorlet{ColourHl3}{TealDrop}
\colorlet{ColourDark}{ICantExpress}
\colorlet{ColourDark2}{Tropiteal}


\usetheme{metropolis}

\setbeamercolor{alerted text}{%
  fg=bazelGreen
}

\setbeamercolor{example text}{%
  fg=mLightBrown
}

\setbeamercolor{frametitle}{%
  use=normal text,
  fg=normal text.bg,
  bg=bazelGreen
}

\lstset{%
  basicstyle=\ttfamily\lst@ifdisplaystyle\fontsize{9pt}{9pt}\selectfont\fi,
  keywordstyle=\color{sgreen}, identifierstyle=\color{sblue}, 
  commentstyle=\color{sbase1}, stringstyle=\color{sorange},
  numberstyle=\color{sviolet}, showstringspaces=false,  
  breaklines=true,
  tabsize=5,
}

\lstalias[]{gnumake}[gnu]{make}





\author{Jonas Rylund Glesaaen}
\title[%
Why are there no good build systems? An introduction to Bazel]{%
Why are there no good build systems? \\ An introduction to Bazel}
\date{June 1st 2018}
\institute{SA2C Tech Chat}

\begin{document}

\maketitle

\begin{frame}{Why do we need build systems?}

  \begin{itemize} \setlength\itemsep{1em}
    \item Essential for big projects
    \item Lowers build times
    \item Incentivises building often
    \begin{itemize}
      \item Makes finding bugs easier
      \item Lowers barrier to testing everything always
    \end{itemize}
  \item Allows for automation (CI/CD etc)
  \item A good system allows you to get on with what you are actually there for
  \end{itemize}

\end{frame}

\begin{frame}{What do I want from a build system?}

  \begin{itemize}
    \item Modular
      \begin{itemize}
        \item Encourages reuse
        \item Simplifies dependencies
        \item I just really like folders
      \end{itemize}
    \item Easy to use
    \item Easy to maintain
    \item Easy to configure
    \item Portable
    \item Extendible to new programming languages
    \item Handles external dependencies cleanly
    \item Lends itself to CI/CD
  \end{itemize}

\end{frame}

\begin{frame}[fragile]
  \frametitle{Let us talk about make}

  It is \alert{\textbf{the}} standard for *NIX based development.

  \vspace{3mm}

  Basically shell + dependencies, the rest you have to make yourself;
  I have too many copies of this snippet:

  \vspace{5mm}

  \begin{lstlisting}[language={[gnu]make}]
SRCS := $(wildcard src/*.cpp)
DEPS := $(SRCS:src/%.cpp=.deps/%.d)

$(DEPS): .deps/%.d: src/%.cpp | .deps
  @$(CXX) $(CXXFLAGS) -MM -MP -MT \
    bin/$(patsubst src/%.cpp,%,$<) -MF $@ $<

-include $(DEPS)
  \end{lstlisting}


\end{frame}

\begin{frame}{Make and recursion}

  \begin{center}
    {\Large\textit{Recursive Make Considered Harmful}} \\
    Miller, P.A. (1998), AUUGN Journal
    % Australian Unix systems User Group
  \end{center}

\end{frame}

\begin{frame}{Make and recursion}

  \begin{itemize}
    \item Hard to get build order right \\
      \emph{Requires a lot of tweaking}
    \item Inter-directory dependencies very difficult \\
      \emph{Builds too much or too little}
    \item Dependencies dropped to decrease build times \\
      \emph{Users need to run clean periodically}
  \end{itemize}

  The \emph{RMCH} paper claims that make does \textbf{\alert{not}} have a design
  flaw, I disagree.

\end{frame}

\begin{frame}{Makefile generators}

  Quite a few "solutions" are simply makefile generators

  \begin{itemize}
    \item autoconf
    \item non-recursive make
    \item CMake
    \item ninja
    \item ...
  \end{itemize}

  We can do better than that!

\end{frame}

\begin{frame}[plain]
  \begin{center}
    {\Large Introducing {\bfseries\color{bazelGreen}Bazel}} \\[1cm]
    \includegraphics[width=0.5\textwidth]{images/bazel-icon.pdf}
  \end{center}
\end{frame}

\begin{frame}{Bazel design principles}

  \begin{itemize} \setlength\itemsep{1em}
    \item Config files are easy to read
    \item \alert{Bazel} is fast and reliable
    \item \alert{Bazel} scales
    \item \alert{Bazel} is extensible
  \end{itemize}

  \BazelNavIcon{}
\end{frame}

\begin{frame}{Bazel concepts}

  \begin{alertblock}{Workspace} \\[5pt]
    The folder that contains your project.\\
    May or may not reference external projects.
  \end{alertblock}

  \vspace{2mm}

  \begin{alertblock}{Packages} \\[5pt]
    Any directory in your workspace that has a {\ttfamily BUILD} file in it.\\
    Contains related files and dependence specifications.
  \end{alertblock}

  \vspace{2mm}

  \begin{alertblock}{Targets} \\[5pt]
    Files and rules the package provides.\\
    Every target has a label.
  \end{alertblock}

  \BazelNavIcon{}
\end{frame}

\begin{frame}{Internal dependencies}

  \alert{Bazel} sandboxes its builds by default.

  \vspace{5mm}

  You therefore have complete control over what is external and what is internal
  for your package.

  \vspace{5mm}

  Also means you need to be in control over package dependencies.

  \vspace{5mm}

  Let us look at Example 01

  \BazelNavIcon{}
\end{frame}

\begin{frame}{External dependencies}

  External as in external to the {\ttfamily WORKSPACE}.

  \vspace{5mm}

  For example {\color{bazelGreen}\ttfamily local\_repository} and
  {\color{bazelGreen}\ttfamily new\_http\_archive}.

  \vspace{5mm}

  Interfaces well with other \alert{bazel} repos, well enough with non-\alert{bazel} ones.

  \vspace{5mm}

  Let us look at Examples 02, 03, \& 04

  \BazelNavIcon{}
\end{frame}

\begin{frame}{The toolchain}

  {
    \only<2->{\color{black!50}}
    So, I wanted to compile an MPI project...
  }

  \vspace{0.5cm}

  \uncover<2->{
    \begin{center}
      \textit{\large This is a fabulously difficult project that causes hardened
      engineers to stare blankly at screens in defeat.}
    \end{center}
  }

  \vspace{0.5cm}

  \uncover<3->{
    Basically to change any one aspect of the compilation process of \alert{bazel}
    you have to replace the entire thing. Let's see an example...
  }

  \BazelNavIcon{}
\end{frame}

\begin{frame}{To sum it up}

  \begin{itemize}
    \item Modular ({\color{bazelGreen}Yes})
    \item Easy to use ({\color{bazelGreen}Yes})
    \item Easy to maintain ({\color{bazelGreen}Yes})
    \item Easy to configure ({\color{mLightBrown}No})
    \item Portable (sort of)
    \item Extendible to new programming languages ({\color{mLightBrown}Not easily})
    \item Handles external dependencies cleanly ({\color{bazelGreen}Yes})
    \item Lends itself to CI/CD ({\color{bazelGreen}Yes})
  \end{itemize}

\end{frame}

\begin{frame}{To sum it up}

  I am not necessarily sold, but as an article by\\
  Windmill Engineering puts it:

  \hfill

  {\large Bazel is the worst build system, except for all the others}

\end{frame}

\end{document}
